% Shared preamble for the theory notes (loaded by theory/main.tex and report/main.tex)
\usepackage[utf8]{inputenc}
\usepackage[T1]{fontenc}
\usepackage{amsmath,amssymb,amsthm,mathtools}
\usepackage{bm}
\usepackage{hyperref}
\usepackage{booktabs}
\usepackage{enumitem}
\usepackage{xcolor}
\usepackage{graphicx}
\usepackage[margin=1in]{geometry}

% ---- theorem environments ----
\newtheorem{theorem}{Theorem}[section]
\newtheorem{proposition}[theorem]{Proposition}
\newtheorem{lemma}[theorem]{Lemma}
\newtheorem{corollary}[theorem]{Corollary}
\newtheorem{definition}[theorem]{Definition}
\newtheorem{assumption}[theorem]{Assumption}
\theoremstyle{remark}
\newtheorem{remark}[theorem]{Remark}
\newtheorem{example}[theorem]{Example}

% ---- notation (matches Bach, Learning Theory from First Principles, Ch. 5) ----
\newcommand{\R}{\mathbb{R}}
\newcommand{\E}{\mathbb{E}}
\newcommand{\Prob}{\mathbb{P}}
\newcommand{\norm}[1]{\left\|#1\right\|}
\newcommand{\inner}[2]{\left\langle #1, #2 \right\rangle}
\newcommand{\abs}[1]{\left|#1\right|}
\newcommand{\etastar}{\eta_\ast}          % minimizer, as in the book
\newcommand{\thetat}{\theta_t}
\newcommand{\grad}{\nabla}
\newcommand{\fip}{f'_i}                   % f'_i = \nabla f_i, book notation
\newcommand{\Fp}{F'}                      % F' = \nabla F
\newcommand{\muEff}{\mu_{\mathrm{eff}}}
\newcommand{\LEff}{L_{\mathrm{eff}}}
\newcommand{\kappaP}{\kappa_P}
\newcommand{\Ei}{\E_{i(t)}}
\newcommand{\tr}{\operatorname{tr}}
\newcommand{\diag}{\operatorname{diag}}
\newcommand{\argmin}{\operatorname*{arg\,min}}
\newcommand{\lmin}{\lambda_{\min}}
\newcommand{\lmax}{\lambda_{\max}}
\newcommand{\ind}{\mathbf{1}}
\DeclareMathOperator{\Var}{Var}

% Table of contents: room for the appendix numbers (B3, A1, ...) and their subsection numbers (B3.10, ...)
\usepackage{etoolbox}
\makeatletter
\patchcmd{\l@section}{1.5em}{2.4em}{}{}
\renewcommand*\l@subsection{\@dottedtocline{2}{2.4em}{3.0em}}
\renewcommand*\l@subsubsection{\@dottedtocline{3}{5.4em}{3.6em}}
\makeatother
